\chapter{Etické a právne prílohy}
\label{chapter:etika}

\section{Informovaný súhlas zákonného zástupcu}

\begin{quote}\small
\textbf{Informovaný súhlas so zapojením do výskumu}

\textbf{Názov výskumu:} Vplyv rôznych foriem podpory umelou inteligenciou na osvojenie učiva informatiky\\
\textbf{Riešiteľ:} Bc. Mykyta Prykhodko, Fakulta informatiky a informačných technológií STU v~Bratislave\\
\textbf{Školiteľ:} Ing. Miroslav Bahleda, PhD.\\
\textbf{Škola:} \rule{4.5cm}{0.4pt} \quad \textbf{Trieda:} \rule{2cm}{0.4pt}

\textbf{O čo ide.} Vaše dieťa absolvuje počas troch po sebe nasledujúcich týždňov tri bežné hodiny informatiky na témy vyhľadávanie, triedenie a rekurzia. Hodiny vedie jeho vlastný učiteľ podľa bežného tematického plánu. Líšia sa iba tým, akú pomoc má žiak k dispozícii pri samostatnej práci: na jednej hodine pracuje bez digitálneho asistenta, na druhej s bežným asistentom typu ChatGPT a na tretej s asistentom upraveným tak, aby nedával hotové riešenia. \textbf{Každý žiak absolvuje všetky tri formy}, takže nikto nie je o žiadnu z nich ukrátený.

\textbf{Čo budeme zisťovať.} Po každej hodine žiaci napíšu krátky test (20 minút) a vyplnia krátky dotazník. Po troch týždňoch napíšu ešte jeden test, ktorý overí, čo si zapamätali. Na hodinách s asistentom zaznamenávame \textbf{text otázok, ktoré žiak asistentovi napíše}, a časové údaje o práci s úlohami.

\textbf{Ako chránime údaje.} Meno žiaka sa v našich údajoch nikdy nenachádza. Nahrádza ho kód, ktorého prevodný kľúč zostáva na škole a výskumník k nemu nemá prístup. Z textu otázok automaticky odstraňujeme mená, adresy a kontaktné údaje. Odpovede asistenta neukladáme v pôvodnom znení. Zaznamenáva sa výhradne komunikácia s asistentom počas vyučovacej hodiny na školskom notebooku; nič mimo hodiny a nič na súkromnom zariadení. Žiaci vidia na začiatku každej hodiny viditeľné upozornenie, že záznam prebieha. Údaje sú uložené šifrovane na infraštruktúre STU, uchovávajú sa tri roky a potom sa odstraňujú.

\textbf{Čo to znamená pre známku.} Nič. Testy z výskumu sa neznámkujú, výsledky sa neoznamujú učiteľovi po jednotlivých žiakoch a účasť ani neúčasť nemá na hodnotenie žiadny vplyv.

\textbf{Dobrovoľnosť.} Účasť je dobrovoľná. Súhlas môžete kedykoľvek a bez uvedenia dôvodu odvolať; vtedy údaje vášho dieťaťa zo súboru odstránime. Ak súhlas neudelíte, vaše dieťa sa hodín zúčastňuje úplne rovnako ako ostatní, len sa jeho údaje nespracúvajú.

\textbf{Právny základ.} Článok 6 ods. 1 písm. a) a článok 8 nariadenia (EÚ) 2016/679 (GDPR) a zákon č. 18/2018 Z. z. Prevádzkovateľom je Slovenská technická univerzita v Bratislave. Máte právo na prístup k údajom, ich opravu, vymazanie a obmedzenie spracúvania, ako aj právo podať sťažnosť Úradu na ochranu osobných údajov SR.

\textbf{Kontakt:} \rule{6cm}{0.4pt} \quad \textbf{Zodpovedná osoba pre ochranu údajov STU:} \rule{4cm}{0.4pt}

\vspace{6pt}
\begin{tabular}{@{}p{7.0cm}p{5.4cm}@{}}
$\square$~Súhlasím so zapojením môjho dieťaťa do výskumu. & \\
$\square$~Nesúhlasím. & \\[10pt]
Meno a priezvisko zákonného zástupcu: \rule{4.5cm}{0.4pt} & Podpis: \rule{3.5cm}{0.4pt} \\[10pt]
Dátum: \rule{3cm}{0.4pt} & \\
\end{tabular}
\end{quote}

\section{Súhlas žiaka}

Súhlas žiaka je samostatný dokument. Text je formulovaný pre sedemnásťročného čitateľa a je zámerne krátky; dlhý text sa nečíta a súhlas potom nie je informovaný.

\begin{quote}\small
\textbf{Chceme ťa poprosiť o súhlas}

Na troch hodinách informatiky budeme skúmať, ako sa dá najlepšie učiť s pomocou umelej inteligencie. Na jednej hodine budeš pracovať bez asistenta, na druhej s bežným asistentom a na tretej s asistentom, ktorý ti nedá hotové riešenie, ale pomôže ti naň prísť.

\textbf{Čo si od teba budeme zaznamenávať:}
\begin{itemize}[nosep]
    \item výsledky krátkych testov po hodinách a jedného testu o tri týždne,
    \item odpovede v krátkych dotazníkoch,
    \item \textbf{text otázok, ktoré napíšeš asistentovi} počas hodiny,
    \item časové údaje o tom, kedy si čo robil/a s úlohami.
\end{itemize}

\textbf{Čo si zaznamenávať nebudeme:} nič mimo vyučovacej hodiny, nič z tvojho telefónu a nič z iných stránok.

\textbf{Dôležité:} tvoje meno v našich údajoch nebude. Testy sa neznámkujú. Nemá to žiadny vplyv na tvoju známku z informatiky. Kedykoľvek môžeš povedať, že už nechceš, a tvoje údaje vymažeme~--- stačí to povedať učiteľovi a nemusíš vysvetľovať prečo.

\textbf{Prosba na záver.} Aby výsledok niečo znamenal, potrebujeme, aby si pracoval/a tak, ako pracuješ normálne. Nie je to skúška toho, aký si dobrý žiak. Zaujíma nás, ako sa s takým nástrojom naozaj pracuje.

\vspace{4pt}
$\square$~Súhlasím \qquad $\square$~Nesúhlasím \qquad Podpis: \rule{3.5cm}{0.4pt} \qquad Dátum: \rule{2.5cm}{0.4pt}
\end{quote}

\section{Posúdenie vplyvu na ochranu údajov}

Posúdenie podľa článku 35 nariadenia GDPR~\cite{gdpr2016} je povinné, pretože spracúvanie zahŕňa systematické monitorovanie a týka sa detí.

\begin{table}[H]
\centering
\caption{Posúdenie vplyvu na ochranu údajov}
\label{tab:dpia}
\scriptsize
\begin{tabularx}{\textwidth}{@{}l X X@{}}
\toprule
\textbf{Riziko} & \textbf{Opis} & \textbf{Opatrenie} \\
\midrule
Odhalenie identity z voľného textu & Žiak môže do promptu napísať svoje meno alebo osobnú informáciu & Automatické odstraňovanie vzorov (mená zo zoznamu, adresy, rodné čísla, kontakty) pred uložením; ručná kontrola vzorky 5\,\%; vykazovanie úspešnosti odstránenia \\
Opätovná identifikácia zo štýlu písania & Štýl je kvázi-identifikátorom & Korpus promptov sa nezverejňuje; prístup len na základe dohody o mlčanlivosti a súhlasu etickej komisie \\
Zachytenie prevádzky nesúvisiacej s výskumom & Proxy vidí všetku prevádzku prehliadača & Zoznam povolených adries; nezhodná prevádzka sa neukladá ani nelogguje; zoznam je súčasťou archivovanej konfigurácie \\
Zneužitie certifikačnej autority & CA na zariadení umožňuje odpočúvanie & Inštalácia len na školské zariadenia zapojené do štúdie; preukázateľné odstránenie po skončení zberu; súkromný kľúč generovaný lokálne a neopúšťa zariadenie \\
Zber mimo vyučovania & Žiak môže notebook použiť aj inokedy & Proxy beží len počas vyučovacieho bloku; spustenie a zastavenie je v protokole vernosti \\
Nerovnováha moci pri súhlase & Žiak môže vnímať účasť ako povinnú & Súhlas vyberá tretia osoba, nie vyučujúci; neúčasť nie je viditeľná pre spolužiakov; žiak zostáva v triede a pracuje rovnako \\
Použitie údajov na hodnotenie žiaka & Učiteľ by mohol chcieť vidieť výsledky & Individuálne výsledky sa učiteľovi neposkytujú; škola dostáva iba agregovanú správu na úrovni triedy pri minimálnom počte 10 žiakov \\
Prenos mimo EHP & Volania modelu smerujú k poskytovateľovi & Použije sa poskytovateľ so spracúvaním v EÚ; ak to nie je možné, doplní sa informácia do súhlasu a zmluvné doložky; prompty sa poskytovateľovi odosielajú bez identifikátorov \\
Strata alebo únik údajov & Notebook sa môže stratiť & Šifrovanie disku; denný prenos na infraštruktúru STU; lokálne kópie sa po prenose odstraňujú \\
\bottomrule
\end{tabularx}
\end{table}

\paragraph{Záver posúdenia.} Zvyškové riziko je po zavedení opatrení vyhodnotené ako prijateľné. Rozhodujúcimi opatreniami sú zoznam povolených adries (ktorý obmedzuje rozsah spracúvania na nevyhnutné minimum), časové ohraničenie zberu na vyučovací blok a uloženie prevodného kľúča mimo dosahu výskumníka. Posúdenie sa predkladá zodpovednej osobe pre ochranu údajov STU pred začatím zberu.

\section{Posúdenie podľa aktu o umelej inteligencii}

Adaptívna vrstva odhaduje úroveň zvládnutia učiva a na jej základe vyberá úlohy, čo je činnosť spadajúca pod bod 3 prílohy III nariadenia (EÚ) 2024/1689~\cite{euaiact2024} (systémy používané vo vzdelávaní na hodnotenie výsledkov učenia). Na systém sa vzťahuje výnimka pre výskumnú a vývojovú činnosť, pretože nie je uvádzaný na trh, nemá vplyv na hodnotenie žiaka a je nasadený výlučne na účel tejto štúdie.

Napriek tomu dokumentujeme systém tak, ako keby sa naň požiadavky vzťahovali, pretože pri prípadnom prechode do školskej praxe by to bolo nevyhnutné a pretože samotná dokumentácia je súčasťou vedeckej hodnoty práce:

\begin{description}[leftmargin=0.6cm, nosep]
    \item[\textbf{Účel a rozsah.}] Systém poskytuje nápovedy pri ôsmich úlohách z troch tém informatiky pre žiakov stredných škôl. Neslúži na hodnotenie, na rozhodovanie o postupe do vyššieho ročníka ani na dohľad počas skúšky.
    \item[\textbf{Údaje.}] Vstupom sú prompty žiaka a stav zvládnutia piatich čiastkových zručností odvodený z jeho pokusov v rámci danej hodiny. Systém nepracuje so žiadnymi údajmi z minulosti žiaka ani s údajmi mimo predmetu.
    \item[\textbf{Známe obmedzenia.}] Jazykový model môže poskytnúť vecne nesprávnu nápoveď. Vrstva to neodstraňuje, len obmedzuje rozsah. Podiel nesprávnych nápovedí sa meria na vzorke 5\,\% odpovedí a vykazuje sa.
    \item[\textbf{Ľudský dohľad.}] Učiteľ je počas celého bloku prítomný, vidí obrazovky žiakov a môže kedykoľvek zasiahnuť. Žiak môže nápoveď odmietnuť a pokračovať samostatne.
    \item[\textbf{Transparentnosť.}] Žiaci sú informovaní, že komunikujú so systémom umelej inteligencie a že systém je nastavený tak, aby nedával hotové riešenia.
    \item[\textbf{Záznam.}] Všetky vstupy a všetky nápovedy sa zaznamenávajú, takže správanie systému je spätne preskúmateľné.
\end{description}

\section{Schvaľovací postup a jeho časovanie}

\begin{table}[H]
\centering
\caption{Schvaľovacie kroky pred začatím zberu}
\label{tab:schvalenia}
\small
\begin{tabularx}{\textwidth}{@{}l X l@{}}
\toprule
\textbf{Krok} & \textbf{Obsah} & \textbf{Najneskôr} \\
\midrule
Etická komisia STU & dizajn, súhlasy, posúdenie vplyvu & T $-$ 8 týždňov \\
Zodpovedná osoba pre ochranu údajov & posúdenie vplyvu, plán správy údajov & T $-$ 8 týždňov \\
Riaditeľ školy & písomný súhlas so zapojením tried & T $-$ 6 týždňov \\
Rodičovská rada & informovanie, nie súhlas & T $-$ 5 týždňov \\
Zákonní zástupcovia & individuálne súhlasy & T $-$ 3 týždne \\
Žiaci & individuálne súhlasy & T $-$ 1 týždeň \\
Predregistrácia & dizajn a plán analýzy vo verejnom registri & T $-$ 1 týždeň \\
\bottomrule
\end{tabularx}
\vspace{2pt}
\footnotesize T = termín prvého merania (vstupná batéria).
\end{table}

Schvaľovací postup je kritickou cestou celého projektu. Zber je viazaný na školský polrok a jeho zmeškanie znamená odklad o celý rok, preto je v časovom pláne (príloha~\ref{chapter:Work-Schedule}) vedený s rezervou štyroch týždňov.

\section{Postup pri nepredvídaných situáciách}

\begin{description}[leftmargin=0.6cm, nosep]
    \item[\textbf{Žiak napíše asistentovi niečo znepokojujúce.}] Prompty sa počas zberu nečítajú v reálnom čase, takže systém nie je nástrojom dohľadu a nemožno sa naň spoliehať pri ochrane žiaka. Ak sa pri neskoršom kódovaní objaví obsah naznačujúci ohrozenie žiaka, postupuje sa podľa vnútorného predpisu školy a kontaktuje sa školský psychológ. Táto možnosť je uvedená v súhlase.
    \item[\textbf{Žiak odvolá súhlas počas zberu.}] Jeho údaje sa odstránia zo všetkých súborov vrátane už vypočítaných agregátov, ktoré sa prepočítajú. Žiak sa naďalej zúčastňuje hodín.
    \item[\textbf{Technický výpadok počas hodiny.}] Hodina pokračuje podľa tlačenej zálohy pracovného listu. Stratia sa procesné ukazovatele, nie učebné výsledky. Prípad sa zaznamená v protokole vernosti a trieda sa v citlivostnej analýze vylúči.
    \item[\textbf{Únik alebo podozrenie na únik údajov.}] Oznámenie zodpovednej osobe do 24 hodín, posúdenie povinnosti oznámenia dozornému orgánu do 72 hodín podľa článku 33 GDPR, informovanie školy a zákonných zástupcov.
    \item[\textbf{Nepriaznivý priebežný výsledok.}] Priebežné výsledky sa počas zberu neanalyzujú. Dôvodom je, že priebežné nazeranie do dát pri neplánovanom zastavení nafukuje chybu prvého druhu. Výnimkou je kontrola technickej funkčnosti zberu (miera pokrytia, počet zachytených promptov), ktorá sa nedotýka výsledkových premenných.
\end{description}
